% GENERATED FILE. Edit canonical lessons, then run npm run generate:books.
% canonical-source-hash: fnv1a64:468fa80e0ac16c2e
% canonical-lessons: KA-C291-tannage, KA-C291-beccage, KA-C291-vicitra, KA-C291-adbhuta, KA-C291-asaktikara

\chapter{Cool, Warm, Strange, Wonderful, Interesting}
\label{ch:ka-cool-warm-strange-wonderful-interesting}
\hlchaptermodality{291}

\begin{chapteropening}
\textbf{By the end of this chapter:} \emph{I can say cool, warm, strange, wonderful and interesting.}

Complete the last lesson of chapter 291: I can say cool, warm, strange, wonderful and interesting.
\end{chapteropening}

\section[\texorpdfstring{\kn{ತಣ್ಣಗೆ} (taṇṇage)}{ತಣ್ಣಗೆ (taṇṇage)}]{\kn{ತಣ್ಣಗೆ} (taṇṇage) — cool, cold (to the touch)}

\label{lesson:KA-C291-tannage}



\begin{quote}
Before the new one: say the Kannada for easy, then the Kannada for important.
\end{quote}

\subsection*{You'll want to know: \kn{ತಣ್ಣಗೆ}}
\textbf{\kn{ತಣ್ಣಗೆ}} — \emph{taṇṇage} — \textquotedblleft{}cool, cold (to the touch)\textquotedblright{}.

\begin{cousinweb}
\textbf{\kn{ತಣ್ಣಗೆ}} (\emph{taṇṇage}) is native Dravidian, from \emph{taṇ}, \textquotedblleft{}coolness, cold\textquotedblright{}, which DEDR 3045 lists, with the adverb ending \emph{-age}; the entry also has Tamil \emph{taṇ}, Telugu \emph{taniyu} and Malayalam \emph{taṇ}. DEDR 3045 compares Prakrit \emph{taṇṇāya}, \textquotedblleft{}damp\textquotedblright{}, and a reconstructed Indo-Aryan \emph{*ṭhaṇḍha}, \textquotedblleft{}cold\textquotedblright{} (CDIAL 13676).
\end{cousinweb}

\begin{grammarlens}[title={What you've built}]
One more word for this chapter.
\end{grammarlens}

\subsection*{Your turn}
\begin{itemize}
\raggedright
  \item \emph{Say it:} \emph{taṇṇage}
  \item \emph{Say it:} \emph{taṇṇage}, once more
  \item \emph{Say it:} say \emph{taṇṇage}
  \item \emph{Recall:} say the Kannada for easy, then the Kannada for important, then say \emph{taṇṇage} again
\end{itemize}

\begin{tcolorbox}[breakable,colback=teal!4,colframe=teal!35!black,title={Before you move on}]
What does \kn{ತಣ್ಣಗೆ} mean? (\textquotedblleft{}Cool, cold (to the touch)\textquotedblright{}.) Say it once more.
\end{tcolorbox}
\section[\texorpdfstring{\kn{ಬೆಚ್ಚಗೆ} (beccage)}{ಬೆಚ್ಚಗೆ (beccage)}]{\kn{ಬೆಚ್ಚಗೆ} (beccage) — warm}

\label{lesson:KA-C291-beccage}



\begin{quote}
Before the new one: say the Kannada for important, then the Kannada for cool.
\end{quote}

\subsection*{You'll want to know: \kn{ಬೆಚ್ಚಗೆ}}
\textbf{\kn{ಬೆಚ್ಚಗೆ}} — \emph{beccage} — \textquotedblleft{}warm\textquotedblright{}.

\begin{cousinweb}
\textbf{\kn{ಬೆಚ್ಚಗೆ}} (\emph{beccage}) is native Dravidian, DEDR 5517 lists \emph{beccage}, \textquotedblleft{}heat, warmth; warm\textquotedblright{}, itself, under the root for \textquotedblleft{}to burn\textquotedblright{}; the entry also has Malayalam \emph{vēka}.
\end{cousinweb}

\begin{grammarlens}[title={What you've built}]
One more word for this chapter.
\end{grammarlens}

\subsection*{Your turn}
\begin{itemize}
\raggedright
  \item \emph{Say it:} \emph{beccage}
  \item \emph{Say it:} \emph{beccage}, once more
  \item \emph{Say it:} say \emph{beccage}
  \item \emph{Recall:} say the Kannada for important, then the Kannada for cool, then say \emph{beccage} again
\end{itemize}

\begin{tcolorbox}[breakable,colback=teal!4,colframe=teal!35!black,title={Before you move on}]
What does \kn{ಬೆಚ್ಚಗೆ} mean? (\textquotedblleft{}Warm\textquotedblright{}.) Say it once more.
\end{tcolorbox}
\section[\texorpdfstring{\kn{ವಿಚಿತ್ರ} (vicitra)}{ವಿಚಿತ್ರ (vicitra)}]{\kn{ವಿಚಿತ್ರ} (vicitra) — strange, odd}

\label{lesson:KA-C291-vicitra}



\begin{quote}
Before the new one: say the Kannada for cool, then the Kannada for warm.
\end{quote}

\subsection*{You'll want to know: \kn{ವಿಚಿತ್ರ}}
\textbf{\kn{ವಿಚಿತ್ರ}} — \emph{vicitra} — \textquotedblleft{}strange, odd\textquotedblright{}.

\begin{cousinweb}
\textbf{\kn{ವಿಚಿತ್ರ}} (\emph{vicitra}) is Sanskrit \emph{vicitra}, \textquotedblleft{}many-coloured, strange\textquotedblright{}, built on \emph{citra}, \textquotedblleft{}bright, variegated\textquotedblright{} (CDIAL 4803).
\end{cousinweb}

\begin{grammarlens}[title={What you've built}]
One more word for this chapter.
\end{grammarlens}

\subsection*{Your turn}
\begin{itemize}
\raggedright
  \item \emph{Say it:} \emph{vicitra}
  \item \emph{Say it:} \emph{vicitra}, once more
  \item \emph{Say it:} say \emph{vicitra}
  \item \emph{Recall:} say the Kannada for cool, then the Kannada for warm, then say \emph{vicitra} again
\end{itemize}

\begin{tcolorbox}[breakable,colback=teal!4,colframe=teal!35!black,title={Before you move on}]
What does \kn{ವಿಚಿತ್ರ} mean? (\textquotedblleft{}Strange, odd\textquotedblright{}.) Say it once more.
\end{tcolorbox}
\section[\texorpdfstring{\kn{ಅದ್ಭುತ} (adbhuta)}{ಅದ್ಭುತ (adbhuta)}]{\kn{ಅದ್ಭುತ} (adbhuta) — wonderful, amazing}

\label{lesson:KA-C291-adbhuta}



\begin{quote}
Before the new one: say the Kannada for warm, then the Kannada for strange.
\end{quote}

\subsection*{You'll want to know: \kn{ಅದ್ಭುತ}}
\textbf{\kn{ಅದ್ಭುತ}} — \emph{adbhuta} — \textquotedblleft{}wonderful, amazing\textquotedblright{}.

\begin{cousinweb}
\textbf{\kn{ಅದ್ಭುತ}} (\emph{adbhuta}) is Sanskrit \emph{adbhuta}, \textquotedblleft{}wonderful\textquotedblright{}, taken over whole as a learned loan, a tatsama (CDIAL 240).
\end{cousinweb}

\begin{grammarlens}[title={What you've built}]
One more word for this chapter.
\end{grammarlens}

\subsection*{Your turn}
\begin{itemize}
\raggedright
  \item \emph{Say it:} \emph{adbhuta}
  \item \emph{Say it:} \emph{adbhuta}, once more
  \item \emph{Say it:} say \emph{adbhuta}
  \item \emph{Recall:} say the Kannada for warm, then the Kannada for strange, then say \emph{adbhuta} again
\end{itemize}

\begin{tcolorbox}[breakable,colback=teal!4,colframe=teal!35!black,title={Before you move on}]
What does \kn{ಅದ್ಭುತ} mean? (\textquotedblleft{}Wonderful, amazing\textquotedblright{}.) Say it once more.
\end{tcolorbox}
\section[\texorpdfstring{\kn{ಆಸಕ್ತಿಕರ} (āsaktikara)}{ಆಸಕ್ತಿಕರ (āsaktikara)}]{\kn{ಆಸಕ್ತಿಕರ} (āsaktikara) — interesting}

\label{lesson:KA-C291-asaktikara}



\begin{quote}
Before the new one: say the Kannada for strange, then the Kannada for wonderful.
\end{quote}

\subsection*{You'll want to know: \kn{ಆಸಕ್ತಿಕರ}}
\textbf{\kn{ಆಸಕ್ತಿಕರ}} — \emph{āsaktikara} — \textquotedblleft{}interesting\textquotedblright{}.

\begin{cousinweb}
\textbf{\kn{ಆಸಕ್ತಿಕರ}} (\emph{āsaktikara}) is Sanskrit \emph{āsakti}, \textquotedblleft{}attachment, interest\textquotedblright{}, from \emph{āsakta}, \textquotedblleft{}attached\textquotedblright{} (CDIAL 1477), and \emph{kara}, \textquotedblleft{}causing\textquotedblright{} (CDIAL 2779): \textquotedblleft{}causing interest\textquotedblright{}.
\end{cousinweb}

\begin{grammarlens}[title={What you've built}]
One more word for this chapter.
\end{grammarlens}

\subsection*{Your turn}
\begin{itemize}
\raggedright
  \item \emph{Say it:} \emph{āsaktikara}
  \item \emph{Say it:} \emph{āsaktikara}, once more
  \item \emph{Say it:} say \emph{āsaktikara}
  \item \emph{Recall:} say the Kannada for strange, then the Kannada for wonderful, then say \emph{āsaktikara} again
\end{itemize}

\begin{tcolorbox}[breakable,colback=teal!4,colframe=teal!35!black,title={Before you move on}]
What does \kn{ಆಸಕ್ತಿಕರ} mean? (\textquotedblleft{}Interesting\textquotedblright{}.) Say it once more.
\end{tcolorbox}
